% LTeX: language=it_IT

\documentclass[main.tex]{subfiles}
\graphicspath{{../Figures/Secondi/}}

\begin{document}
    \chapter{Secondi piatti}

    \newpage
    \begin{recipe}
        [%
            preparationtime = {\unit[45]{min}},
            portion = {\unit[1]{Porzione}},
            source = {Lyon 16/10/2021},
        ]
        {Roesti di patate e zucca al Camembert}

        \graph{
            big = RoestiPatateZucca.JPG
        }

        \ingredients{
            1/2 & Cipolla \\
            \unit[150]{g} & Zucca \\
            1 & patata \\
            \unit[50]{g} & Camembert \\
            1 & Uovo \\
            \unit[1]{cucchiaio} & Olio EVO \\
            \unit[1]{bicchiere} & Acqua \\
            \unit[1]{pizzico} & Sale \\
            \unit[1]{pizzico} & Pepe nero \\
            \unit[1]{pizzico} & Zucchero \\
        }

        \preparation{
            \begin{enumerate}
            \item Mondate la cipolla e tagliatela a pezzettini piccoli. Scaldate l'olio in una padella antiaderente e soffriggete.
            \item Tagliate a cubetti piccoli la zucca e le patate ed aggiungetele alla padella. Fate saltare per un minuto circa, poi aggiungete gradualmente l'acqua. Il quantitativo di acqua dipende da quanto è alto il fuoco e se cucinate con o senza coperchio. L'importante è non annegare il contenuto ma neanche bruciarlo. Aggiungere sale, pepe e zucchero e mescolare.
            \item Quando le patate e la zucca si sono ammorbidite fate evaporare l'acqua in eccesso (lasciate comunque che il composto sia umido) e poi aggiungete il Camembert. Fatelo sciogliere e mescolate il tutto, mantecando.
            \item Aggiungete l'uovo e mescolate per amalgamare.
            \item Smettete di mescolare e aspettate che il composto assuma una consistenza pseudo-solida.
            \item Servite accompagnando con un insalata verde semplice (e.g. Valerianella).
            \end{enumerate}
        }
    \end{recipe}

    \newpage
    \begin{recipe}
        [%
            preparationtime = {\unit[40]{min}},
            portion = {\unit[3]{Porzioni}},
            source = {Copenhagen 17/09/2021},
        ]
        {Dahl di lenticchie e fagioli neri}

        \graph{
            big = Dahl.jpeg
        }

        \ingredients{
            2 & Cipolle \\
            \unit[3]{Spicchi} & Aglio \\
            \unit[400]{g} & Lenticchie in scatola \\
            \unit[100]{g} & Fagioli neri in scatola \\
            \unit[60]{g} & Uva nera fresca \\
            \unit[3]{cucchiai} & Olio EVO \\
            \unit[1/2]{bicchiere} & Acqua \\
            \unit[100]{g} & Passata di pomodoro \\
            \unit[1]{pizzico} & Sale \\
            \unit[1]{pizzico} & Pepe nero \\
            2 & Peperoncini \\
            \unit[2]{cucchiai} & Miele \\
        }

        \preparation{
            \begin{enumerate}
            \item Mondate la cipolla e tagliatela grossolanamente. Pulite e schiacciate gli spicchi d'aglio. Mettete il tutto a soffriggere in padella.
            \item Quando la cipolla è dorata, aggiungete le lenticchie e i fagioli con un po' della loro acqua. Aggiungete anche i due pepeoncini, mescolate e fate cuocere per 10 minuti con il coperchio.
            \item Aggiungente alla padella la passata di pomodoro, il sale, il pepe e il miele. Mescolate e continuate la cottura.
            \item Lavate e tagliate in 4 gli acini d'uva e aggiungeteli alla padella. Mescolate e cuocete ancora per qualche minuto, questa volta senza coperchio, in modo da asciugare il composto.
            \item Quando il Dahl ha raggiunto una consistenza cremosa, spegnete il fuoco e servite.
            \end{enumerate}
        }
    \end{recipe}

    \newpage
    \begin{recipe}
        [%
            preparationtime = {\unit[1]{ora}},
            portion = {\unit[2]{Porzioni}},
            source = {Idefix 17/08/2022},
        ]
        {Pane Frattau}

        \graph{
            big = PaneFrattau.jpeg
        }

        \ingredients{
            \unit[1]{confezione} & Pana Carrasau \\
            \unit[250]{g} & Sugo di pomodoro \\
            2 & Uova \\
            \unit[100]{g} & Pecorino sardo \\
            \unit[100]{g} & Parmigiano \\
            \unit[1]{L} & Acqua \\
            1 & Dado vegetale
        }

        \preparation{
            \begin{enumerate}
            \item Se preferite, invece del sugo di pomodoro già pronto, potete farlo voi al momento. Se usate il sugo pronto, scaldatelo in un pentolino.
            \item In una pentola più larga dei dischi di pane carrasau, mettete a bollire l'acqua e shoglietevi il dado vegetale.
            \item Grattuggiate il pecorino e il Parmigiano.
            \item Preparate due piatti fondi, con mezzo mestolo di sugo ben spalmato.
            \item Uno alla volta, mettete i dischi di pane carrasau nel brodo, pescateli prima che si disfino e poneteli nei piatti, alternando con mezzo mestolo di sugo di pomodoro e una spolverata di formaggio gratuggiato
            \item Quando avete ottenuto il numero desiderato di strati, uno alla volta, versate le uova crude nel brodo per farle in camicia. Pescatele facendo attenzione a non romperle e disponetele in cima alla torre di pane, una per piatto.
            \item Mangiate subito ché si raffredda in fretta.
            \end{enumerate}
        }
    \end{recipe}


\end{document}